\part{Topological Quantum Phenomena}

\chapter{A field-free loop can carry phase}
\label{ch:ab}

\physicalquestion{How can an electron acquire a measurable phase while
traveling only through a region where the electromagnetic field vanishes?}

\modelassumptions{An ideal solenoid confines magnetic flux to an inaccessible
region.  Outside it the particle is coherent, nonrelativistic, and described
by minimal $U(1)$ coupling.  Leakage fields and environmental decoherence are
neglected in the ideal model.}

\section{What does the interference observation demand?}

Two electron paths passing on opposite sides of a shielded solenoid recombine.
The model predicts a relative phase even though $F=dA=0$ along each path.
The gauge-invariant datum is not the value of $A$ at a point but the holonomy
\[
\operatorname{Hol}_A(\gamma)
=\exp\!\left(\frac{iq}{\hbar}\int_\gamma A\right)
\]
around the closed loop obtained from the two paths.

In the field-free region $M$, the potential is a closed one-form.  A gauge
change $A\mapsto A+d\lambda$ changes its representative but not its de Rham
class $[A]\in H^1_{\mathrm{dR}}(M)$.

\section{Why is holonomy topological in the flat region?}

\begin{theorem}[Cohomology--homology dependence]
For a closed one-form $A$ and a closed loop $\gamma$, the integral
$\int_\gamma A$ depends only on $[A]\in H^1_{\mathrm{dR}}(M)$ and
$[\gamma]\in H_1(M;\R)$.  Consequently the Aharonov--Bohm phase is invariant
under gauge change and homology-preserving deformation of the path.
\end{theorem}

\begin{proof}
Replacing $A$ by $A+d\lambda$ changes the integral by
$\int_\gamma d\lambda=0$.  If $\gamma-\gamma'=\partial C$, then Stokes'
theorem from Chapter~\ref{ch:forms} gives
\[
\int_\gamma A-\int_{\gamma'}A=\int_CdA=0.
\]
Both dependencies therefore descend to the stated classes.
\end{proof}

\section{What phase does the ideal solenoid predict?}

For $M=\R^3$ minus the solenoid axis, let
\[
A=\frac{\Phi}{2\pi}d\theta.
\]
A once-winding loop has $\int_\gamma A=\Phi$, so the relative phase is
$q\Phi/\hbar$ modulo $2\pi$.  Shielded-flux electron interferometry measures
the displacement of fringes as flux varies.  The experiment tests this
holonomy prediction while separately controlling leakage fields; it does not
prove de Rham theory \cite{aharonov-bohm,tonomura}.

\section{What further contact should be proved?}

\begin{exercise}[Winding computes holonomy]
For a loop of winding number $n$ around the removed axis, prove
$\int_\gamma d\theta=2\pi n$ and compute its phase.
\end{exercise}

\begin{exercise}[Flux periodicity]
Show that all interference probabilities are unchanged when
$q\Phi/\hbar$ changes by $2\pi k$.  Identify the corresponding flux quantum.
\end{exercise}


\chapter{Curvature produces stable characteristic classes}
\label{ch:chern}

\physicalquestion{Which global quantity detects a twisted family of quantum
states even though every small patch admits a smooth frame?}

\modelassumptions{The parameter space is a compact smooth manifold and the
relevant eigenspaces have constant finite rank.  Connections are smooth and
the structure group is unitary.}

\section{How can local curvature define a global class?}

For a complex vector bundle $E\to X$ with connection $\nabla$, let
$F_\nabla$ be its curvature.  The total Chern form is defined by
\[
\det\!\left(1+\frac{i}{2\pi}F_\nabla\right)
=1+c_1(E,\nabla)+c_2(E,\nabla)+\cdots.
\]
The forms are assembled from invariant polynomials, so they agree under
changes of local frame.

Stable equivalence declares $E$ and $F$ equivalent when
$E\oplus\underline\C^r\cong F\oplus\underline\C^s$ after balancing ranks.
The Grothendieck group of bundles under direct sum is $K^0(X)$.

\section{Why do Chern classes ignore the chosen connection?}

\begin{theorem}[Chern--Weil independence]
Each Chern form is closed, and its de Rham class is independent of the
connection.  Direct sum makes the Chern character
\[
\operatorname{ch}(E)=\left[\tr\exp\!\left(\frac{iF_\nabla}{2\pi}\right)\right]
\]
descend to a homomorphism
$\operatorname{ch}:K^0(X)\to H^{\mathrm{even}}_{\mathrm{dR}}(X)$.
\end{theorem}

\begin{proof}
The Bianchi identity from Chapter~\ref{ch:gauge} and invariance of trace imply
$d\,\tr(F_\nabla^k)=0$.  For a path of connections $\nabla_t$ with derivative
$a_t$, differentiating under the trace gives
\[
\frac d{dt}\tr(F_t^k)=k\,d\,\tr(a_t\wedge F_t^{k-1}).
\]
Integrating in $t$ shows that two choices differ by an exact form.  The
exponential trace is additive on direct sums, so it respects the Grothendieck
relation and descends to $K^0(X)$.
\end{proof}

\section{What invariant results for a surface?}

For a complex bundle over a closed oriented surface,
\[
C_1(E)=\frac{i}{2\pi}\int_X\tr(F_\nabla)
\]
is an integer.  It is unchanged by a smooth deformation of the connection or
by adding trivial bundles.  Curvature may shift from place to place, but its
total Chern number cannot vary continuously.  The integer can jump only when
the underlying bundle ceases to vary continuously within the assumed class.

\section{What further contact should be proved?}

\begin{exercise}[A nonzero Chern number obstructs a global frame]
Prove that a bundle admitting a global frame is trivial and has vanishing
positive-degree Chern classes.  Conclude that $C_1(E)\ne0$ obstructs a global
smooth eigenbasis.
\end{exercise}

\begin{exercise}[Stable equivalence preserves the Chern character]
Show directly that adding a trivial bundle changes only the degree-zero rank
term of the Chern character.
\end{exercise}


\chapter{The Hall integer is a bundle invariant}
\label{ch:hall}

\physicalquestion{Why does Hall conductance remain pinned to integer
multiples of $e^2/h$ while microscopic Hamiltonian parameters vary?}

\modelassumptions{The crystal is noninteracting and translation invariant,
the Fermi energy lies in a spectral gap, temperature is negligible compared
with that gap, and the occupied bands vary smoothly over the Brillouin torus.
The TKNN linear-response formula is taken as the physical bridge from the
model to conductance.}

\section{How does a gapped Hamiltonian produce topology?}

Let $H:k\mapsto H(k)$ be a smooth family of finite-dimensional self-adjoint
operators over the Brillouin torus $\mathbb T^2$.  A uniform gap at zero gives
the occupied spectral projection
\[
P(k)=\frac{1}{2\pi i}\int_\Gamma(z-H(k))^{-1}\,dz,
\]
where $\Gamma$ surrounds the negative spectrum.  The images of $P(k)$ form
the Bloch bundle $E_{\mathrm{occ}}\to\mathbb T^2$.  Adding inert filled or
empty flat bands changes the bundle by a trivial summand, so the stable phase
is a K-theory class.

\section{Why can the Hall integer change only at a gap closing?}

\begin{theorem}[Gap-preserving invariance]
A smooth gap-preserving homotopy $H_s(k)$ gives isomorphic endpoint occupied
bundles.  Hence
\[
\mathrm{Ch}_1(P)
=\frac{i}{2\pi}\int_{\mathbb T^2}
\tr\bigl(P\,dP\wedge dP\bigr)
\]
is an integer independent of $s$ and of added trivial bands.
\end{theorem}

\begin{proof}
Uniform compactness of $[0,1]\times\mathbb T^2$ keeps a common contour inside
the gap, so the Riesz formula produces a smooth projection on the product.
Its image is a vector bundle whose restrictions at $s=0$ and $s=1$ are
isomorphic because an interval is contractible.  The displayed form
represents the first Chern class by Chapter~\ref{ch:chern}; its integral is
therefore constant and integral.  A trivial summand has zero curvature and
does not change the pairing.
\end{proof}

\section{How does topology return to conductance?}

The TKNN formula predicts
\[
\sigma_{xy}=\mathrm{Ch}_1(P)\frac{e^2}{h}.
\]
This is the physical bridge from the bundle invariant to the observed
response \cite{tknn,von-klitzing}.
Thus continuous changes of hopping amplitudes cannot move the ideal
conductance away from its plateau while the gap and model assumptions remain
valid.  A jump requires a point where the occupied projection is no longer
defined continuously---typically a gap closing.  Disorder and interactions
require broader formulations, but the invariant-versus-gap logic survives in
operator-algebraic versions.

\section{What further contact should be proved?}

\begin{exercise}[Berry curvature equals projection curvature]
For an isolated occupied line with local unit vector $\psi(k)$, compare
$i\langle\psi,d\psi\rangle$ with the connection induced by $P$.  Prove that
their curvatures yield the projection formula above.
\end{exercise}

\begin{exercise}[A two-band transition]
For $H(k)=d(k)\cdot\sigma$ with $d(k)\ne0$, show that the occupied Chern
number is the degree of $d/|d|:\mathbb T^2\to S^2$.  Identify why the degree
can change only where $d=0$.
\end{exercise}


\chapter{Fredholm index counts protected zero modes}
\label{ch:index}

\physicalquestion{Why is the imbalance of two kinds of zero modes stable
even though each kernel dimension can jump under perturbation?}

\modelassumptions{Operators act between complex Hilbert spaces, have closed
range, and have finite-dimensional kernel and cokernel.  Perturbations are
bounded and small in norm, or compact when explicitly stated.}

\section{Which difference survives perturbation?}

An operator $T:H_0\to H_1$ is Fredholm when its kernel and cokernel are finite
dimensional and its range is closed.  Its index is
\[
\operatorname{ind}T=\dim\ker T-\dim\operatorname{coker}T.
\]
The compact operators $\mathcal K(H)$ form an ideal in the bounded operators,
and the Calkin algebra is $\mathcal Q(H)=\mathcal B(H)/\mathcal K(H)$.

\section{Why is the index topologically stable?}

\begin{theorem}[Atkinson stability]
$T$ is Fredholm exactly when its image in the Calkin algebra is invertible.
The Fredholm index is locally constant and is unchanged by compact
perturbations.
\end{theorem}

\begin{proof}
If $T$ is Fredholm, choose inverses on complements of its finite-dimensional
kernel and cokernel; extending by zero gives a parametrix $S$ for which
$ST-1$ and $TS-1$ have finite rank.  Conversely a parametrix modulo compacts
implies finite-dimensional kernel and cokernel and closed range.  Thus
Fredholm operators are the inverse image of the invertible group of
$\mathcal Q(H)$, an open set.  Along a sufficiently small norm path the
kernel-cokernel difference cannot change, as follows from the parametrix and
finite-dimensional reduction.  Compact perturbations have the same Calkin
image and can be joined through $T+sK$, so the index is unchanged.
\end{proof}

\section{How does index meet geometry and physics?}

An elliptic differential operator on a closed manifold has a Fredholm
extension.  The Atiyah--Singer theorem identifies its analytic index with a
topological K-theory pairing of its principal symbol \cite{atiyah-singer}.
In gauge backgrounds,
this converts characteristic charge into chiral zero-mode imbalance.  The
deep theorem is used here as a bridge; the stability proved above already
explains why admissible perturbations cannot change the imbalance.

For the circle operator $D=d/d\theta$ on periodic complex functions, kernel
and cokernel are both the constants, so $\operatorname{ind}D=0$, agreeing
with the Euler characteristic of $S^1$.  This elementary case displays the
analytic-topological equality without claiming a proof of the general index
theorem.

\section{What further contact should be proved?}

\begin{exercise}[The unilateral shift has index minus one]
On $\ell^2(\mathbb N)$ let $S(e_n)=e_{n+1}$.  Compute the kernels and
cokernels of $S$ and $S^*$, then verify index stability after a finite-rank
perturbation.
\end{exercise}

\begin{exercise}[Part IV synthesis]
For the Aharonov--Bohm phase, Hall integer, and Fredholm index, identify the
structured input, the deformation class, the numerical pairing, and the
hypothesis whose failure permits a jump.  Express the first as a
cohomology--homology pairing, the second as a Chern pairing of a K-class, and
the third as the boundary invariant of an invertible Calkin class.
\end{exercise}

\parttransition{Topology explains quantized sectors but not how maps behave
after arbitrary matrix amplification, how nonlocal correlations are bounded,
or how infinite observable algebras contain physical weak limits.  Those
questions require noncommutative analysis.}
